% ============================================================
% FloodRSCT: SIGSPATIAL 2026 Application Paper — front matter
% Title, author, abstract, Introduction (§1–§2)
% ============================================================
%
% Build notes:
%   - Targets ACM acmart class with the sigconf option (SIGSPATIAL standard).
%   - During early drafting you may prefer the NeurIPS-style article preamble
%     used in combine_georsct.tex; the alternative preamble is included
%     commented-out at the bottom of this header block. Switch by toggling
%     the two blocks.
%   - Bib entries for \cite{martin2026rsct} and \cite{martin2026georsct}
%     are provided in floodrsct.bib (companion file).
%
% ============================================================
% PREAMBLE (ACM sigconf)
% ============================================================

\documentclass[sigconf,nonacm=true,review,authoryear]{acmart}

% Suppress ACM reference format block during review drafting.
\settopmatter{printacmref=false}
\renewcommand\footnotetextcopyrightpermission[1]{}
\pagestyle{plain}

% Use natbib author-year citation commands (\citep, \citet) consistent
% with the GeoRSCT draft.
\setcitestyle{authoryear,round}

\usepackage[utf8]{inputenc}
\usepackage[T1]{fontenc}
\usepackage{booktabs}
\usepackage{amsmath}
\usepackage{amssymb}
\usepackage{bm}
\usepackage{microtype}
\usepackage{graphicx}
\usepackage{xcolor}
\usepackage{algorithm}
\usepackage{algpseudocode}
\usepackage{subcaption}
\usepackage{url}

% Safety net for tight two-column lines: allow modest extra stretch
% before declaring overfull. Less aggressive than \sloppy.
\setlength{\emergencystretch}{1em}

% ============================================================
% ALTERNATIVE PREAMBLE (article + NeurIPS-style, for early drafting)
% Uncomment the block below and comment out the acmart block above
% if you want to compile with the same setup as combine_georsct.tex.
% ============================================================
%
% \documentclass{article}
% \usepackage[final]{neurips_2026}
% \usepackage[utf8]{inputenc}
% \usepackage[T1]{fontenc}
% \usepackage{hyperref}
% \usepackage{url}
% \usepackage{booktabs}
% \usepackage{amsfonts}
% \usepackage{amsmath}
% \usepackage{bm}
% \usepackage{microtype}
% \usepackage{graphicx}
% \usepackage{xcolor}
% \usepackage{algorithm}
% \usepackage{algpseudocode}
% \usepackage{subcaption}
% \usepackage[round,authoryear]{natbib}    % provides \citep / \citet
%
% ============================================================

% ACM conference metadata (placeholder; replace at camera-ready)
\acmConference[SIGSPATIAL '26]{The 34th ACM SIGSPATIAL International
  Conference on Advances in Geographic Information Systems}{November
  3--6, 2026}{Riverside, CA, USA}
\acmYear{2026}
\copyrightyear{2026}

\title{FloodRSCT: Certificate-Based Diagnostics for Crisis-Ready Flood Forecasting [Applications]}
% Note: per SIGSPATIAL 2026 Applications Track CFP, papers must include
% the [Applications] suffix in the submission title; it is removed at
% camera-ready. Single-blind submission, so author block remains visible.

\author{Rudolph A. Martin}
\affiliation{%
  \institution{Next Shift Consulting LLC}
  \city{Madison}
  \state{Wisconsin}
  \country{USA}
}
\email{rudy@nextshiftconsulting.com}

\begin{document}

% ============================================================
% ABSTRACT
% ============================================================

\begin{abstract}
Flood crisis systems increasingly use spatial prediction models to prioritize warnings, evacuations, inspections, and resource allocation. Scalar flood-risk scores, however, do not indicate whether a warning is reliable under spatial shift, missing sensors, infrastructure fragility, scale mismatch, or proxy-driven prediction. We introduce \textbf{FloodRSCT}, a certificate-based diagnostic framework that separates flood risk from forecast readiness. For each location, FloodRSCT combines predicted risk with signal quality, representation--solver compatibility, perturbation stability, task difficulty, and geospatial stress evidence, then maps the result to operational actions: \textbf{Trust}, \textbf{Review}, \textbf{Escalate}, or \textbf{Suppress}. The framework is composed of four separable layers---forecast, certificate, retrieval, and crisis action---each independently auditable. The certificate machinery is inherited from Representation--Solver Compatibility Theory~\citep{martin2026rsct} and was empirically validated on a non-flood spatial benchmark in GeoRSCT~\citep{martin2026georsct}; FloodRSCT adapts and extends it for flood crisis decision support. Across five U.S. flood-crisis archetypes spanning Houston, New Orleans, Riverside--Coachella, Southwest Florida, and NYC/NJ, we show how certificate diagnostics reveal failures hidden by scalar accuracy and support auditable crisis decision workflows. We further evaluate operator-facing rationales as a separate layer, distinguishing forecast accuracy, certificate reliability, retrieval grounding, and explanation acceptability. The framework is demonstrated in \textbf{Floodcaster.com}, an operator-facing system that loads 72-hour flood scenarios, produces location-level certificates, generates an action queue, and supports auditable drill-down into modal evidence.
\end{abstract}

\keywords{flood forecasting, crisis decision support, spatial prediction,
  representation--solver compatibility, model risk, RAG grounding,
  emergency management}

\maketitle

% ============================================================
% SECTION 1: INTRODUCTION
% ============================================================
\section{Introduction}
\label{sec:intro}

\subsection{Crisis Motivation}
\label{sec:intro-motivation}

A 72-hour flood forecast ranks several neighborhoods as high risk. Emergency managers cannot act on a risk score alone: they need to know which warnings are reliable, which need analyst review, which require escalation despite uncertainty, and which should be suppressed because the model is outside its reliable operating zone. This problem is especially acute during flood crises because missed warnings can be catastrophic, but false confidence can also misdirect scarce field teams, evacuation support, and infrastructure inspection capacity~\citep{morss2008}. The operational question is not which model scores highest in aggregate, but which warning is decision-ready for this place, at this time, under this scenario.

Most flood-model evaluation reports scalar accuracy: AUC, RMSE, $R^2$, $F_1$, calibration error, or hit rate. These scores are useful but incomplete. A model can perform well on average while failing in specific neighborhoods, relying on socioeconomic proxies instead of flood evidence, breaking under small rainfall perturbations, or overfitting spatial leakage~\citep{roberts2017cv,ploton2020spatial}. Scalar accuracy answers ``how well does this model rank average performance,'' not ``should I trust this warning for this ZCTA in the next 72 hours.''

We distinguish \textbf{forecast risk} from \textbf{forecast readiness}. Forecast risk estimates whether flooding may occur. Forecast readiness asks whether the warning is reliable enough to support action. The two can diverge: a high-risk warning may require escalation because its evidence is fragile, while a moderate-risk warning may be trusted because its flood-modal signal is stable and geographically compatible. FloodRSCT is the diagnostic layer that makes this distinction operational.

\subsection{What FloodRSCT Contributes}
\label{sec:intro-contributes}

FloodRSCT applies Representation--Solver Compatibility Theory (RSCT)~\citep{martin2026rsct} to flood crisis forecasting. RSCT introduces a per-instance certificate that decomposes evaluation evidence into signal quality $\alpha$, representation--solver compatibility $\kappa$, representational turbulence $\sigma$, and a task-difficulty bound $\mathrm{TRF}$, all defined over an $(R, S_{\text{sup}}, N)$ simplex. GeoRSCT~\citep{martin2026georsct} validated this machinery on a 27-task CONUS-scale geospatial benchmark, demonstrating that scalar performance can be substrate-conditioned (architecture sensitivity is categorical by target type, not continuous in task difficulty). FloodRSCT inherits these foundations and extends them with flood-specific geometry diagnostics, a geospatial stress-test protocol, and an action-mapping gate that produces operational decisions rather than scalar priority scores. Certificate definitions are reproduced in Appendix~\ref{app:cert-arch} for self-containment; underlying invariance proofs and cross-substrate validation are cited rather than reproduced.

We evaluate FloodRSCT across five U.S. flood-crisis archetypes selected to stress distinct hydrometeorological regimes and operational failure modes: Houston/Harris County urban pluvial flooding, New Orleans protected-basin infrastructure flooding, Riverside--Coachella desert flash flooding, Southwest Florida storm surge, and NYC/NJ dense urban cloudburst flooding. The five cases are not minor variants of one phenomenon; they exercise different decision geometries (drainage networks, levees and pumps, canyon runoff, surge exposure, basement vulnerability, scale mismatch). The five scenarios are therefore evidence of \emph{generality} of one certificate framework rather than five hand-built systems: one certifier, one Gate~4 grounding contract, five evidence registries, five compiled narrative templates.

\paragraph{Four-layer separation.}
FloodRSCT is composed of four separable layers (Figure~\ref{fig:four-layers}): a \textbf{forecast layer} that estimates flood risk; a \textbf{certificate layer} that diagnoses forecast readiness from model evidence and geospatial stress tests; an \textbf{evidence retrieval (RAG) layer} that supplies authority-ranked sources for grounding the certificate's gate decisions and for generating operator rationales; and a \textbf{crisis-action layer} that maps certificate output to Trust, Review, Escalate, or Suppress. The certificate is \emph{not} computed from retrieved evidence; retrieved evidence is used to test whether the forecast's grounding gate fires defensibly and to supply citations in the operator rationale. This separation is what makes the system auditable: the forecast can be wrong, the certificate can be reliable, the retrieved evidence can be authoritative, and the rationale can still be unacceptable---each is evaluated against its own contract.

\begin{figure*}[t]
\centering
\includegraphics[width=\textwidth]{figures/floodrsct_four_layer_separation.pdf}
\caption{The four-layer separation. Forecast estimates risk; certificate diagnoses readiness; retrieval supplies authority-ranked grounding evidence (R/S/N-annotated, snapshot-immutable, Tier-1-precedence); crisis-action gate maps certificate output to Trust/Review/Escalate/Suppress. Layers are evaluated independently: forecast accuracy, certificate reliability, retrieval grounding, and rationale acceptability are four separate failure surfaces.}
\Description{A horizontal pipeline diagram showing four colored boxes
  in sequence. From left to right: a blue Forecast box (spatial model
  emits per-location risk score), an orange Certificate box (the
  alpha-kappa-sigma-TRF tuple plus geometry flags), a green
  Retrieval box (Tier 1/2/3 evidence with R/S-sup/N annotations), and
  a red Crisis Action box (fixed-order gate emitting Trust, Review,
  Escalate, or Suppress). Solid arrows connect the boxes left to
  right, labeled with the handoffs: risk score, gate query, and
  evidence plus R/S/N annotations. Below the pipeline, a gray band
  shows four evaluation surfaces (forecast accuracy, certificate
  reliability, retrieval grounding, rationale acceptability), each
  connected to its corresponding pipeline layer by a dashed arrow.
  The band is labeled with the principle that each layer is judged
  against its own contract.}
\label{fig:four-layers}
\end{figure*}

\paragraph{The Floodcaster.com demonstration system.}
We implement the FloodRSCT framework in \textbf{Floodcaster.com}, an operator-facing system that loads a 72-hour flood scenario, produces location-level certificates, generates an action queue, supports drill-down into modal evidence (county $\rightarrow$ ZCTA $\rightarrow$ modal source $\rightarrow$ feature), runs scenario-specific stress-test recipes, and produces a plain-language operator brief with an audit-trace record per action. \emph{The contribution of this paper is FloodRSCT}; Floodcaster.com is the system instantiation that operationalizes the method and the artifact through which the empirical results are produced. Floodcaster.com does not replace official warnings, hydrologic authorities, or field validation (\S\ref{sec:governance}); it provides diagnostic support for interpreting flood-model outputs.

\subsection{Contributions}
\label{sec:intro-contributions}

This paper makes five contributions:

\begin{enumerate}
\item \textbf{A flood forecast readiness certificate} that separates predicted risk from reliability. For each location and 72-hour scenario, the certificate combines $\alpha$ (signal quality), $\kappa$ (representation--solver compatibility), $\sigma$ (perturbation stability), $\mathrm{TRF}$ (task difficulty), and geometry flags into an auditable diagnostic record, then maps it via a fixed-order gate to a four-valued operational action vocabulary \{Trust, Review, Escalate, Suppress\} (\S\ref{sec:method}).

\item \textbf{A geospatial stress-test protocol} for missing sensors, spatial perturbation, infrastructure failure, scale mismatch, and topology-sensitive residuals along drainage networks, levee/pump segments, canyon routing, surge zones, and dense urban infrastructure. Stress tests are implemented as deterministic recipes over a single ablation engine---not as duplicated per-scenario endpoints---so that the same compute path is exercised across all five scenarios (\S\ref{sec:method}, Appendix~\ref{app:stress-recipes}).

\item \textbf{An auditable evidence-retrieval architecture} with a snapshot-immutable Postgres backbone, three-tier authority sources (regulatory, operational SOP, technical), a deterministic scenario router with allowed/required/forbidden evidence-family registries per gate question, and an $(R, S_{\text{sup}}, N)$ annotation contract that makes the retrieval layer's output a first-class evidence object in the same simplex the certifier operates on. Generic retrieve-then-read RAG is explicitly deprioritized: it lacks the authority tier, provenance chain, gate-question-relative $(R, S_{\text{sup}}, N)$ annotation, and snapshot determinism required for certificate-grade grounding (\S\ref{sec:method}, Appendix~\ref{app:rag-arch}).

\item \textbf{Evaluation across five hydrometeorological flood regimes} rather than a single local dashboard. Each scenario exercises distinct failure modes---drainage and proxy reliance (Houston), protected-basin infrastructure fragility (New Orleans), terrain/topology instability (Riverside--Coachella), surge tail-event uncertainty (Southwest Florida), and dense urban scale mismatch (NYC/NJ)---and the same certificate machinery is shown to produce scenario-appropriate action distributions and dominant gate-failure signatures (\S\ref{sec:scenarios}, \S\ref{sec:results}).

\item \textbf{A rationale-evaluation layer} that assesses operator-facing explanations independently of forecast accuracy and certificate reliability. Adapting the self-rationalization out-of-distribution evaluation discipline of~\citet{rocha2024selfrat}, we evaluate whether each rationale justifies its assigned action, cites supported evidence, and avoids overclaim, with a shortcoming taxonomy (insufficient justification, unsupported evidence, vagueness, overconfidence, proxy/demographic overclaim, ignored reliability warning, ignored high-risk condition, none). Rationale acceptability is reported as a fourth, independent failure surface alongside forecast accuracy, certificate reliability, and retrieval grounding (\S\ref{sec:method}, \S\ref{sec:results}).
\end{enumerate}

\paragraph{Relationship to RSCT and GeoRSCT.}
This is an application paper. The certificate framework (RSCT~\citep{martin2026rsct}) and its cross-substrate empirical validation (GeoRSCT~\citep{martin2026georsct}) are inherited; we cite them for foundational definitions, invariance proofs, the $\mathrm{TRF}$ block-bootstrap procedure, and the cross-task benchmark evidence that the framework's empirical claims generalize beyond any one substrate. The contributions of FloodRSCT are the flood-specific instantiation of the gate order, the geospatial stress-test protocol, the snapshot-immutable RAG grounding architecture, the five-scenario evaluation, the rationale evaluation discipline, and the Floodcaster.com operator-facing instantiation. The appendix is structured to preserve the self-contained property: definitions that the reader must \emph{operate on} are reproduced; results that \emph{certify the foundations} are cited (Appendix~\ref{app:cert-arch} states this discipline explicitly).

\paragraph{Paper organization.}
Section~\ref{sec:related} reviews flood forecasting, crisis decision support, spatial validation, geometric spatial learning, and rationale evaluation. Section~\ref{sec:problem} formalizes the flood crisis scenario and the certificate's inputs and outputs. Section~\ref{sec:method} presents the FloodRSCT method, including the fixed-order action gate and the evidence-retrieval architecture. Section~\ref{sec:scenarios} introduces the five hydrometeorological regimes. Section~\ref{sec:experiments} states the six research questions and the experimental design. Section~\ref{sec:results} reports results across all six research questions. Section~\ref{sec:demo} describes Floodcaster.com. Section~\ref{sec:governance} states governance and non-claims. Section~\ref{sec:limitations} discusses limitations, and Section~\ref{sec:conclusion} concludes.

% ============================================================
% (Subsequent sections — Related Work, Problem Formulation,
%  Method, Scenarios, Experiments, Results, Demo, Governance,
%  Limitations, Conclusion — to be drafted in subsequent passes.)
% ============================================================

\bibliographystyle{ACM-Reference-Format}
\bibliography{floodrsct}

\end{document}
